\makeabstract
{
Developing an Optical Microscopy Method to Detect Microplastics in Environmental Samples
}
{
Rowan Bixler (1), Mohammadjavad Hajirezaei (2), Dr. Ryan Poling-Skutvi, PhD (2)
}
{
\textsuperscript{1}Department of Physics, Amherst College, Amherst, MA\\
\textsuperscript{2}Department of Chemical Engineering, University of Rhode Island, Kingston, RI
}
{
Microplastic contamination in coastal environments poses a growing threat to ecosystems and human health, yet current detection methods for these microscopic particles are often limited by long processing times, high costs, limited availability, and low signal detection for complex samples. This project aims to develop an optical microscopy technique using Differential Dynamic Microscopy (DDM) to identify and characterize microplastics in environmental water samples. DDM analyzes light intensity fluctuations to extract information about particle dynamics and properties. This approach could enable effective characterization and differentiation of microplastic types, enhancing tracking, monitoring, and mitigation efforts.
}
{
Microplastic samples composed of polystyrene and polymethyl methacrylate (PMMA) were suspended in water and imaged using brightfield microscopy. Videos of particles were recorded at varying focal planes to capture particle dynamics across sample depth. These videos were processed using a DDM code to extract quantitative parameters, including amplitude as a function of the wavevector (q) and diffusivity. Additional image analysis was conducted to examine and quantify ring patterns produced by particle light scattering at varying axial positions.
}
{
Analysis revealed that diffusivity increased with height from the substrate bottom, eventually reaching a plateau near the middle of the sample—consistent with expected Brownian motion behavior. The amplitude function, A(q), exhibited local minima corresponding to the particles’ effective length scale. When plotted against axial position (z), the amplitude peak followed a Gaussian distribution. Fitting this profile enabled the determination of the particle’s refractive index, yielding a value (\textasciitilde{}1.53) that aligns closely with the known range for polystyrene (1.55–1.59).
}
{
These findings lay the foundation for an efficient, accessible, and label-free method of microplastic identification and provide insight into how DDM can be adapted for complex, real-world samples. Future work will focus on quantifying and refining measurement error, as well as expanding the analysis to include additional polymer types and mixed samples to broaden applicability and validate trends.
}

\newpage
